% ===== PRÉAMBULE CANONIQUE — à coller VERBATIM en tête de CHAQUE .tex =====
\documentclass[11pt,a4paper]{article}
\usepackage[utf8]{inputenc}\usepackage[T1]{fontenc}\usepackage{lmodern}
\usepackage[french]{babel}
\usepackage{amsmath,amssymb,mathtools}\usepackage{icomma}
\usepackage{siunitx}\sisetup{locale=FR}  % \qty{}{}, \si{}, \ang{}
\usepackage{xcolor}\usepackage{colortbl}
\usepackage{tikz}\usepackage{pgfplots}\pgfplotsset{compat=1.18}
\usepackage[siunitx,europeanresistors]{circuitikz} % circuits (élec.)
\usepackage[most]{tcolorbox}\usepackage{enumitem}\usepackage{array,booktabs}
\usepackage[a4paper,margin=2.2cm]{geometry}
\usetikzlibrary{arrows.meta,calc,angles,quotes,positioning,patterns,%
  backgrounds,shapes.geometric,decorations.pathmorphing,decorations.markings,intersections}
\definecolor{primary}{RGB}{222,49,99}\definecolor{primarydark}{RGB}{150,22,55}
\definecolor{defcol}{RGB}{0,114,178}\definecolor{methcol}{RGB}{52,92,148}
\definecolor{excol}{RGB}{0,135,90}\definecolor{attncol}{RGB}{200,40,55}
\tcbset{boxbase/.style={enhanced,breakable,boxrule=0.9pt,arc=2.5pt,
  left=3mm,right=3mm,top=2.3mm,bottom=2.3mm,fonttitle=\bfseries,coltitle=white}}
% --- boîtes TUTORIEL (noms EXACTS, ne pas renommer) ---
\newtcolorbox{cle}[1][]{boxbase,colback=defcol!6,colframe=defcol,
  title={\textsc{À retenir}\ifx\relax#1\relax\relax\else\textmd{ : \itshape #1}\fi}}
\newtcolorbox{methode}[1][]{boxbase,colback=methcol!7,colframe=methcol,sharp corners,
  title={\textsc{Méthode}\ifx\relax#1\relax\relax\else\textmd{ --- \itshape #1}\fi}}
\newtcolorbox{exemplebox}[1][]{boxbase,colback=excol!5,colframe=excol,
  title={\textsc{Exemple}\ifx\relax#1\relax\relax\else\textmd{ : \itshape #1}\fi}}
\newtcolorbox{piege}[1][]{boxbase,colback=attncol!6,colframe=attncol,
  title={\textsc{Piège}\ifx\relax#1\relax\relax\else\textmd{ : \itshape #1}\fi}}
\newtcolorbox{remarque}[1][]{boxbase,colback=black!4,colframe=black!45,
  title={\textsc{Remarque}\ifx\relax#1\relax\relax\else\textmd{ : \itshape #1}\fi}}
% --- boîtes EXERCICE / CORRIGÉ (noms EXACTS, auto-comptées) ---
\newtcolorbox[auto counter]{exo}[1][]{boxbase,colback=primary!4,colframe=primary,
  title={\textsc{Exercice}~\thetcbcounter\ifx\relax#1\relax\relax\else\textmd{ --- \itshape #1}\fi}}
\newtcolorbox[auto counter]{corrige}[1][]{boxbase,colback=excol!4,colframe=excol,
  title={\textsc{Corrigé}~\thetcbcounter\ifx\relax#1\relax\relax\else\textmd{ --- \itshape #1}\fi}}
\newcommand{\vv}[1]{\vec{#1}}\newcommand{\ds}{\displaystyle}
\newcommand{\R}{\mathbb{R}}\newcommand{\N}{\mathbb{N}}\newcommand{\Z}{\mathbb{Z}}
% =========================================================================
\begin{document}
\usetikzlibrary{babel}
\begin{corrige}[trois lampes : vrai ou faux ?]
Les trois lampes sont identiques. $L_1$ porte tout le courant du circuit ; il se partage ensuite entre
$L_2$ et $L_3$ (loi des nœuds au point $A$) : $I=I_2+I_3$.
\begin{enumerate}[label=\alph*)]
  \item \textbf{Vrai}. $L_2$ et $L_3$ sont en parallèle (même tension) et identiques (même résistance),
        donc parcourues par la même intensité : $I_2=I_3$.
  \item \textbf{Vrai}. Comme $I_2=I_3$, la loi des nœuds donne $I=I_2+I_3=2I_2$ : $L_1$, traversée par
        $I$, porte bien le \emph{double} de $I_2$.
  \item \textbf{Faux}. $L_1$ est en série avec le générateur (aucun nœud entre eux) : les deux
        ampèremètres mesurent la même intensité $I$.
  \item \textbf{Vrai}. $I=I_2+I_3=2I_2$, donc $I_2=I_3=\dfrac{I}{2}=\dfrac{0{,}6}{2}=\qty{0.3}{\ampere}$.
\end{enumerate}
\end{corrige}

\begin{corrige}[deux nœuds en cascade]
On applique la loi des nœuds successivement.
\begin{itemize}[leftmargin=1.4em,topsep=2pt]
  \item Au nœud $N_1$ : $I=I_a+I_b \Rightarrow I_a=I-I_b=8-3=\qty{5}{\ampere}$.
  \item Au nœud $N_2$ : $I_a+I_c=I_d \Rightarrow I_d=5+2=\qty{7}{\ampere}$.
\end{itemize}
\end{corrige}

\begin{corrige}[comptage d'électrons — style concours]
En $t=\qty{1}{\second}$ : $Q=I\,t=\qty{3.2}{\coulomb}$, d'où
\[ n=\frac{Q}{e}=\frac{3{,}2}{1{,}6\times10^{-19}}=2{,}0\times10^{19}\ \text{électrons/s}. \]
(La section de \qty{4}{\milli\metre\squared} est une donnée \emph{inutile} : le comptage ne dépend que
de $I$ et de $t$.)
\begin{enumerate}[label=\alph*)]
  \item \textbf{Faux} : une mole vaut $6{,}02\times10^{23}\gg2{,}0\times10^{19}$ (on est $\sim3\times10^4$
        fois en dessous d'une mole).
  \item \textbf{Faux} : $2{,}0\times10^{19}>10^{13}$.
  \item \textbf{Vrai} : $2{,}0\times10^{19}>1{,}8\times10^{13}$.
\end{enumerate}
\end{corrige}

\begin{corrige}[l'éclair : un courant intense et bref]
\begin{enumerate}[label=\alph*)]
  \item $t=\qty{2}{\milli\second}=2\times10^{-3}\,\si{\second}$, donc
        \[ I=\frac{Q}{t}=\frac{15}{2\times10^{-3}}=\qty{7500}{\ampere}. \]
  \item $n=\dfrac{Q}{e}=\dfrac{15}{1{,}6\times10^{-19}}=9{,}4\times10^{19}$ électrons.
  \item $\dfrac{7500}{10}=750$ : l'éclair débite une intensité environ \textbf{750 fois} plus grande
        que celle d'un circuit domestique — mais pendant un temps très bref.
\end{enumerate}
\end{corrige}

\end{document}
